\documentclass[11pt,a4paper]{article}

\usepackage[margin=23mm,headheight=15pt]{geometry}
\usepackage[T1]{fontenc}
\usepackage[utf8]{inputenc}
\usepackage{lmodern}
\usepackage[scaled=0.92]{helvet}
\usepackage{microtype}
\usepackage{parskip}
\usepackage{amsmath,amssymb,bm}
\usepackage{booktabs,tabularx,array,longtable}
\usepackage{graphicx}
\usepackage{float}
\usepackage{xcolor}
\usepackage{enumitem}
\usepackage{caption}
\usepackage{fancyhdr}
\usepackage{hyperref}
\usepackage[nameinlink,noabbrev]{cleveref}

\definecolor{ink}{HTML}{17324D}
\definecolor{muted}{HTML}{67727E}
\definecolor{rulegray}{HTML}{D9DEE3}
\definecolor{statusgreen}{HTML}{27864B}
\definecolor{statusgreenlight}{HTML}{E7F4EA}
\definecolor{statusamber}{HTML}{C48319}
\definecolor{statusamberlight}{HTML}{FFF2D8}
\definecolor{statusred}{HTML}{B5443C}
\definecolor{statusredlight}{HTML}{F8E5E3}
\definecolor{accent}{HTML}{168A85}
\definecolor{linkblue}{HTML}{2C6E9B}

\hypersetup{
  colorlinks=true,
  linkcolor=ink,
  citecolor=linkblue,
  urlcolor=linkblue,
  pdftitle={NEGFPy 1.0 bulk phonon and SMRTA validation report},
  pdfauthor={NEGFPy validation project}
}
\captionsetup{
  font=small,
  labelfont={bf,color=ink},
  textfont={color=ink},
  margin=8pt
}
\setlist[itemize]{leftmargin=1.55em,itemsep=2pt,topsep=3pt}
\setlist[enumerate]{leftmargin=1.75em,itemsep=2pt,topsep=3pt}
\setlength{\emergencystretch}{3em}
\renewcommand{\arraystretch}{1.16}

\makeatletter
\renewcommand\@seccntformat[1]{\csname the#1\endcsname\hspace{0.65em}}
\renewcommand\section{\@startsection{section}{1}{\z@}%
  {-2.6ex \@plus -0.8ex \@minus -.2ex}%
  {1.2ex \@plus .2ex}%
  {\normalfont\Large\sffamily\bfseries\color{ink}}}
\renewcommand\subsection{\@startsection{subsection}{2}{\z@}%
  {-2.0ex \@plus -0.6ex \@minus -.2ex}%
  {0.8ex \@plus .2ex}%
  {\normalfont\large\sffamily\bfseries\color{ink}}}
\renewcommand\subsubsection{\@startsection{subsubsection}{3}{\z@}%
  {-1.6ex \@plus -0.4ex \@minus -.2ex}%
  {0.6ex \@plus .2ex}%
  {\normalfont\normalsize\sffamily\bfseries\color{ink}}}
\makeatother

\input{generated/evidence_macros.tex}
\input{generated/harmonic_table.tex}
\input{generated/mesh_table.tex}

\pagestyle{fancy}
\fancyhf{}
\fancyhead[L]{\small\sffamily\color{muted}NEGFPy bulk phonon/SMRTA validation}
\fancyhead[R]{\small\sffamily\color{muted}\CandidateShortCommit}
\fancyfoot[L]{\small\sffamily\color{muted}Version 1.0.0 candidate}
\fancyfoot[R]{\small\sffamily\color{muted}\thepage}
\renewcommand{\headrulewidth}{0.4pt}
\renewcommand{\footrulewidth}{0pt}

\newcommand{\vect}[1]{\bm{#1}}
\newcommand{\q}{\vect{q}}
\newcommand{\kvec}{\vect{k}}
\newcommand{\code}[1]{\texttt{#1}}
\newcommand{\scopewarning}[1]{%
  \begin{quote}\small\color{muted}\sffamily #1\end{quote}%
}

\begin{document}
\hypersetup{pageanchor=false}

\begin{titlepage}
\thispagestyle{empty}
\vspace*{9mm}

{\sffamily\bfseries\color{accent}\large NEGFPy VALIDATION PROGRAM\par}
\vspace{8mm}
{\sffamily\bfseries\color{ink}\fontsize{27}{32}\selectfont
Bulk phonons and intrinsic thermal\par conductivity\par}
\vspace{2mm}
{\sffamily\color{muted}\fontsize{16}{20}\selectfont
Band structure, density of states, three-phonon relaxation,\par
and SMRTA transport\par}

\vspace{11mm}
\noindent{\color{rulegray}\rule{\textwidth}{0.7pt}}
\vspace{8mm}

\noindent\fcolorbox{\EvidenceStatusColor}{\EvidenceStatusBackground}{%
  \parbox{\dimexpr\textwidth-2\fboxsep-2\fboxrule\relax}{%
    \vspace{2mm}
    {\sffamily\bfseries\color{\EvidenceStatusColor}\large
      \EvidenceStatusTitle\par}
    \vspace{1mm}
    {\small \EvidenceStatusMessage\par}
    \vspace{2mm}
  }%
}

\vspace{7mm}
\small
\begin{tabularx}{\textwidth}{>{\sffamily\bfseries\color{ink}}p{0.28\textwidth} X}
Candidate revision & \CandidateCommit \\
Git tree & \CandidateTree \\
Reference implementation & phono3py 4.4.0 \\
Reference dataset & Si-PBEsol example pinned at commit
\texttt{302ad205452e\ldots} \\
Transport approximation & Intrinsic three-phonon single-mode relaxation-time approximation (SMRTA) \\
Candidate release gates & \textbf{\CandidateDecisionStatus} \\
Certification authority & \texttt{release-completion-receipt.json} (sole) \\
Authoritative decision & \textbf{\AuthoritativeDecisionStatus} \\
Physical mesh convergence & \textbf{false} \\
Report input time & \ReportGeneratedUTC \\
\end{tabularx}
\normalsize

\vfill
\noindent{\small\color{muted}
The numerical values on the $3^3$, $5^3$, and $11^3$ meshes are
implementation-validation points evaluated at fixed settings. They are not a
mesh- or smearing-converged prediction of silicon thermal conductivity. This
report does not validate anharmonic heat transport through an open junction.}
\end{titlepage}

\hypersetup{pageanchor=true}
\pagenumbering{roman}
\tableofcontents
\clearpage
\pagenumbering{arabic}

\section{Executive result and scope}

\ifreportcertified
The frozen commit-last completion receipt, acting as the sole certification
authority, records that the exact version 1.0.0 candidate release gates passed.
The bound validation JSON and evidence sidecar remain provisional with
\code{release\_certified=false}, as required by the acyclic publication
contract. This statement certifies the bulk-observable and SMRTA implementation
only; \emph{physical\_mesh\_convergence\_certified remains false}.
\else
This is a provisional rendering. Either the sole completion receipt does not
certify the candidate, or this is a legacy development input without the frozen
commit-last authority. Recorded same-input comparisons may still be useful,
but no release tag or silicon-property claim follows from this document.
\fi

The 1.0 scope is a bulk-crystal workflow. Given mapped second- and third-order
interatomic force constants (FC2 and FC3), it supplies:

\begin{itemize}
  \item harmonic frequencies and eigenvectors on arbitrary paths or complete
  uniform meshes;
  \item phonon band structures, total density of states (DOS), and analytic
  group velocities;
  \item on-shell intrinsic three-phonon linewidths and population lifetimes;
  \item mode heat capacities, mean free paths, tensorial SMRTA conductivity,
  and cumulative conductivity; and
  \item immutable numerical results carrying units, input identities,
  approximation names, and provenance metadata.
\end{itemize}

\scopewarning{The 1.0 conductivity is a bulk phonon-Boltzmann result. It is
not evaluated from a left-device-right junction, and its reservoirs do not
contain FC3. Open-device anharmonic self-energies and their conservation laws
belong to the planned 2.0 development line.}

\section{Formulation and unit contract}

\subsection{Harmonic spectrum and band structure}

With mass-normalized reciprocal-space dynamical matrix $D(\q)$, the harmonic
problem is
\begin{equation}
  D(\q)\vect{e}_{\q s}=\omega_{\q s}^{2}\vect{e}_{\q s},
  \qquad
  \vect{e}_{\q s}^{\dagger}\vect{e}_{\q s'}=\delta_{ss'}.
\end{equation}
The public frequency is the cyclic frequency
$\nu_{\q s}=\omega_{\q s}/(2\pi)$ in THz. A band structure samples the same
eigenproblem along a declared piecewise-linear reciprocal-space path. It is
never reconstructed by sorting an unrelated uniform mesh.

For real-space FC2, the release preflight checks masses, atom mappings,
reciprocal pairs, translational and rotational residuals, random-$\q$
Hermiticity, $D(-\q)=D(\q)^{*}$, and sampled stability before harmonic
observables can be called release-eligible.

\subsection{Density of states}

For a complete uniform mesh with $N_q$ points and $n_b$ branches, the total
DOS per primitive cell is
\begin{equation}
  g(\nu)=\frac{1}{N_q}\sum_{\q s}\delta(\nu-\nu_{\q s}),
  \qquad
  \int g(\nu)\,d\nu=n_b=3N_{\mathrm{atom}}.
\end{equation}
The implementation supports an explicitly declared broadening and refuses to
treat a band path as a DOS sampling mesh. A frequency window that truncates
the normalization is rejected unless the caller opts into a separately named
truncated-window result.

\subsection{Group velocity}

The Cartesian group velocity is
\begin{equation}
  \vect{v}_{\q s}=\nabla_{\kvec}\omega_{\q s},
\end{equation}
evaluated analytically from derivatives of $D(\q)$. Within a degenerate
subspace the derivative matrices are diagonalized before velocities are
reported. The archived unit is THz\,\AA, with
$1\ \mathrm{THz}\,\text{\AA}=100\ \mathrm{m\,s^{-1}}$.

\subsection{Three-phonon linewidth and population lifetime}

The FC3 calculation evaluates the standard lowest-order three-phonon
self-energy using energy-conserving decay and difference channels, Bose
occupation factors, and crystal-momentum conservation. Normal and Umklapp
contributions sum to the total on-shell linewidth. The production input
$\gamma_{\q s}$ follows the phono3py convention: cyclic-frequency HWHM in THz.
Therefore
\begin{equation}
  \Gamma^{\mathrm{angular}}_{\mathrm{FWHM},\q s}
  =4\pi\gamma_{\q s}\quad [\mathrm{ps}^{-1}],
  \qquad
  \tau_{\q s}=\frac{1}{4\pi\gamma_{\q s}}\quad [\mathrm{ps}].
  \label{eq:lifetime}
\end{equation}
This factor of $4\pi$ is a locked unit identity, not a fitted convention.
Modes below the declared acoustic cutoff or with unresolved nonpositive rates
are masked explicitly rather than assigned an arbitrary lifetime.

\subsection{Heat capacity, mean free path, and SMRTA conductivity}

For $x=h\nu/(k_BT)$, the harmonic mode heat capacity is
\begin{equation}
  C_{\q s}=k_B\frac{x^2e^x}{(e^x-1)^2}.
\end{equation}
Small-$x$, large-$x$, and zero-temperature limits are evaluated without
unstable direct exponentials. The mode mean-free-path vector and the bulk
conductivity tensor are
\begin{align}
  \vect{\Lambda}_{\q s} &= \vect{v}_{\q s}\tau_{\q s},\\
  \kappa_{ij}^{\mathrm{SMRTA}}(T)
  &=\frac{1}{N_q V}\sum_{\q s}
    C_{\q s}(T)v_{\q s,i}v_{\q s,j}\tau_{\q s}(T),
  \label{eq:kappa}
\end{align}
where $V$ is the primitive-cell volume. The stored per-mode tensors must sum
back to the total tensor. Raw band-index decompositions are bookkeeping only:
within a degenerate manifold, only the complete subspace sum is gauge
invariant.

\section{Approximation hierarchy}

\begin{table}[htbp]
\centering
\caption{Physics included in version 1.0 and its explicit boundaries.}
\small
\begin{tabularx}{\textwidth}{>{\bfseries\raggedright\arraybackslash}p{0.21\textwidth} >{\raggedright\arraybackslash}p{0.18\textwidth} >{\raggedright\arraybackslash}X}
\toprule
Layer & 1.0 status & Meaning and limitation \\
\midrule
Harmonic lattice dynamics & Included & FC2 bands, DOS, eigenvectors, and analytic group velocities. \\
Intrinsic three-phonon scattering & Included & Lowest-order on-shell FC3 rates with normal and Umklapp channels; FC3 is not placed in a lead or open device. \\
SMRTA & Included & Each mode relaxes independently. Off-diagonal mode coupling and the iterative linearized BTE correction are excluded. \\
Isotopes and boundaries & Excluded & No mass-disorder, finite-size, grain-boundary, or phenomenological boundary rate is silently added. \\
Higher anharmonicity & Excluded & No quartic renormalization, four-phonon scattering, thermal expansion, or tadpole shift. \\
Nonphononic channels & Excluded & No electron-phonon, defect, magnetic, or radiative scattering. \\
Coherent/interband transport & Excluded & No Wigner interband term; no off-diagonal density-matrix conductivity. \\
Anharmonic NEGF junction & Planned 2.0 & Cubic device self-energies, Dyson/Keldysh iteration, current conservation, and transverse-momentum coupling. \\
\bottomrule
\end{tabularx}
\end{table}

SMRTA is mature and useful for rapid screening, mode-resolved diagnosis, and
an independent bulk scale for later device calculations. It is not universally
quantitative. In crystals dominated by momentum-conserving normal processes,
the iterative BTE may differ materially because the collision operator is not
diagonal in the harmonic-mode basis \cite{ziman1960,togo2015}.

\section{IFC, mapping, and provenance gates}

The release workflow separates raw input, audit, correction, mass
normalization, and transport assembly. Raw arrays are never silently
symmetrized, remapped, range-truncated, or assigned missing reciprocal partners.
The same crystal mapping and phase convention are shared by FC2 reconstruction,
FC3 Fourier transformation, and mode-space contractions.

\begin{table}[htbp]
\centering
\caption{Required input and provenance checks. ``Current artifact'' refers to
the machine-readable input used by this rendering, not to an inferred state.}
\small
\begin{tabularx}{\textwidth}{>{\bfseries\raggedright\arraybackslash}p{0.25\textwidth} >{\raggedright\arraybackslash}p{0.27\textwidth} >{\raggedright\arraybackslash}X}
\toprule
Gate & Release requirement & Current artifact evidence \\
\midrule
FC2 physical audit & Reciprocity, ASR, rotational moments, Hermiticity,
time reversal, stability, masses, and mappings pass declared tolerances &
\ifreportcertified Direct harmonic metrics and the FC2 schema/physical-audit
decisions are authenticated and passing.\else Numerical audit details must be
embedded in the final 1.0 validation JSON.\fi \\
FC3 identity and audit & Full/compact mapping, permutation/ASR/range audit,
and physical-input SHA bind the linewidth field & \ifreportcertified The sealed
source and archive identities are bound through the release evidence
sidecar.\else The development artifact records pinned FC3 and v0.2 source
hashes; final 1.0 must authenticate the sealed v0.2 manifest.\fi \\
Linewidth identity & Mesh, selected modes, cutoff, smearing, units, and
physics-only problem digest match the RTA problem & Enforced by the result
assembler and HDF5 archive invariants. \\
Source provenance & Exact clean commit, tree, wheel, environment lock, and
input hashes agree & \ifreportcertified The recorded before/after runtime
states are identical, clean, and input-authenticated.\else A dirty development
checkout cannot satisfy this gate.\fi \\
Archive integrity & Atomic regular files, strict schemas, payload SHA256, and
commit-last cross-binding & The provisional JSON and evidence remain false;
the sole receipt binds their hashes, the HDF5 hash, the exact-tree checksum
manifest hash, and the candidate decision before any certified wording. \\
\bottomrule
\end{tabularx}
\end{table}

The exact report JSON, release-evidence sidecar, HDF5, checksum-manifest, and
completion-receipt identities used for a rendering are written to
\path{generated/normalized_report_state.json}. The report generator refuses
symlinked inputs, reconstructs the checksum manifest from the complete artifact
tree, verifies the evidence SHA embedded by both the JSON and HDF5, verifies
the receipt's report/evidence/HDF5/manifest hashes and candidate decision, and
recomputes all per-mode conductivity sums.

\section{Silicon implementation validation}

\subsection{External comparisons}

The pinned silicon comparison uses the phono3py 4.4.0 Si-PBEsol example and
the exact FC2/FC3 artifacts already qualified by the 0.2.0 bulk oracle. The
group-velocity reference uses phono3py's analytic implementation; the heat
capacity and conductivity references are assembled independently from the
same archived frequencies and rates. Small $3^3$ and $5^3$ cases additionally
run phono3py end to end. \ifreportcertified The authenticated $11^3$ harmonic
archive also supplies the fixed $\Gamma$--X--L--$\Gamma$ band path and DOS,
whose external error metrics pass the declared gates.\else A final 1.0 report
also requires externally compared high-symmetry bands and DOS, rather than
only plotting internally generated curves.\fi

\begin{table}[H]
\centering
\caption{Harmonic-observable reporting fields. \ifreportcertified These rows
are authenticated by the locked release evidence.\else Missing rows keep this
rendering uncertified even if other numerical comparisons pass.\fi}
\small
\begin{tabularx}{\textwidth}{>{\bfseries}p{0.30\textwidth} p{0.25\textwidth} X}
\toprule
Observable & Status & Maximum absolute error \\
\midrule
\HarmonicValidationRows
\bottomrule
\end{tabularx}
\end{table}

\begin{table}[H]
\centering
\caption{Same-input silicon implementation validation. Conductivities are in
W\,m$^{-1}$\,K$^{-1}$ and are reported at the single archived temperature.
The values are not a mesh-convergence series.}
\scriptsize
\begin{tabular}{ccccccc}
\toprule
Mesh & $N_q$ & max $|\Delta v|$ & max $|\Delta C|$ & NEGFPy $\kappa_{xx}$ & reference $\kappa_{xx}$ & end to end \\
 & & (THz\,\AA) & (eV K$^{-1}$) & & & \\
\midrule
\MeshValidationRows
\bottomrule
\end{tabular}
\end{table}

The variation among the three archived values is evidence that these points
must \emph{not} be interpreted as a converged silicon property. Their purpose
is code-path equality at fixed inputs.

\begin{figure}[htbp]
  \centering
  \includegraphics[width=\textwidth]{figures/band_and_dos.pdf}
  \caption{\textbf{Harmonic observables.} \ifreportcertified The left panel
  shows the authenticated fixed $\Gamma$--X--L--$\Gamma$ band path; the right
  panel shows the authenticated total DOS. Their independently computed
  comparison errors are reported in the harmonic-validation table.\else A
  final report plots the independently compared high-symmetry band path at
  left. The provisional artifact lacks that dataset, so the panel remains an
  explicit placeholder. At right, the draft plot is a Gaussian DOS computed
  only from the complete $11^3$ uniform-mesh frequencies; it is not promoted
  to an external DOS validation.\fi}
  \label{fig:band-dos}
\end{figure}

\begin{figure}[htbp]
  \centering
  \includegraphics[width=\textwidth]{figures/rates_and_mfp.pdf}
  \caption{\textbf{Three-phonon relaxation and transport lengths on the
  $11^3$ implementation mesh.} Population lifetimes use the exact cyclic-HWHM
  conversion in \cref{eq:lifetime}. Mean free paths combine the same lifetimes
  with analytic group velocities. Points below the acoustic cutoff or without
  a positive resolved rate are omitted. For readability the mean-free-path
  panel also omits displayed values below $10^{-6}$\,nm; all numerical values
  remain in the HDF5 artifact.}
  \label{fig:rates}
\end{figure}

\begin{figure}[htbp]
  \centering
  \includegraphics[width=\textwidth]{figures/conductivity_and_cumulative.pdf}
  \caption{\textbf{SMRTA conductivity diagnostics.} The left panel deliberately
  uses a logarithmic axis to show the nonconverged variation among the three
  implementation meshes. The right panel accumulates the nonnegative trace
  contribution of $11^3$ modes by mean-free-path threshold. For display,
  contributions at or below $10^{-6}$\,nm are collected at that lower bound;
  the accumulated conductivity is unchanged. Neither panel is evidence of
  physical mesh convergence.}
  \label{fig:kappa}
\end{figure}

\clearpage
\section{Numerical invariants and release gates}

The 1.0 result objects and archives enforce cross-field identities in addition
to pointwise comparisons:
\begin{itemize}
  \item $\Gamma_{\mathrm{FWHM}}^{\mathrm{angular}}=4\pi\gamma$ and
  $\tau\Gamma_{\mathrm{FWHM}}^{\mathrm{angular}}=1$ on every resolved mode;
  \item $\vect{\Lambda}=\vect{v}\tau$ under the stored THz\,\AA/ps convention;
  \item included-mode, positive-rate, stability, cutoff, and q-point masks are
  identical across the linewidth and conductivity inputs;
  \item every per-mode tensor sums to the archived total tensor, which is
  symmetric within the strict field tolerances;
  \item compact and full FC2/FC3 representations share physical mappings and
  reproduce gauge-invariant quantities; and
  \item serial and task-level MPI bulk-linewidth fields are merged in canonical
  target-$\q$ order without floating-point MPI reductions.
\end{itemize}

Final release acceptance requires exact Python 3.10 and 3.12 test records, a
fresh-wheel import, clean commit/tree/environment provenance, authenticated
reuse of the 0.2.0 linewidth oracle, external band/DOS comparison, the complete
$3^3/5^3/11^3$ matrix, strict archives, and an independently checked rendered
PDF. A post-processing or report failure can resume from immutable completed
scientific artifacts; it must not trigger an unnecessary FC3 rerun.

\section{Version 1.0 boundary and version 2.0 plan}

\begin{table}[htbp]
\centering
\caption{Scientific boundary between the first bulk release and the planned
anharmonic-NEGF development line.}
\small
\begin{tabularx}{\textwidth}{>{\bfseries\raggedright\arraybackslash}p{0.24\textwidth} >{\raggedright\arraybackslash}X >{\raggedright\arraybackslash}X}
\toprule
Question & Version 1.0 & Planned version 2.0 \\
\midrule
Geometry & Infinite periodic bulk crystal & Left lead--finite device--right lead, with optional transverse periodicity \\
Anharmonicity & On-shell FC3 rates used in a bulk SMRTA collision model & Frequency-dependent cubic NEGF self-energy in the device \\
Leads & Not applicable to bulk BTE & Harmonic in the first material-device stage; anharmonic reservoirs require a separate formalism \\
Primary output & $\nu$, DOS, $\vect{v}$, $\gamma$, $\tau$, $\vect{\Lambda}$, and $\bm{\kappa}$ & Spectral heat current, interacting conductance, collision power, and spatial conservation \\
Core validation & phono3py equality and bulk invariants & Dense/sparse SCBA equality, acoustic contact limits, current conservation, and $q$-coupled convergence \\
\bottomrule
\end{tabularx}
\end{table}

The first 2.0 vertical slice should therefore map the already validated silicon
FC3 into a padded $N_q=1$ interaction region, prove dense/sparse equality, and
demonstrate contact-current and internal-collision-power conservation. Only
after that gate should transverse-momentum-coupled SCBA and silicon-film or
interface comparisons begin.

\clearpage
\section{Reproducibility record}

This report is generated from one provisional validation JSON, one bound HDF5
results archive, one provisional release-evidence sidecar, their exact-tree
checksum manifest, and the commit-last completion receipt. The completion
receipt is the sole certification authority; neither provisional input may
carry a true release flag.
The report JSON SHA256 is \ValidationJSONSHA; the release-evidence SHA256 is
\ReleaseEvidenceSHA; the whole-file HDF5 SHA256 is \ResultsHDFSHA; the checksum
manifest SHA256 is \ChecksumManifestSHA; and the completion-receipt SHA256 is
\CompletionReceiptSHA. The generator creates figures, LaTeX evidence macros,
a normalized report state, and a PDF checksum list. It
does not call FC2, FC3, phono3py, or any transport solver. Figure generation is
deterministic for fixed numerical inputs; the PDF build uses a fixed
\code{SOURCE\_DATE\_EPOCH}, UTC, and disabled shell escape.

The source report deliberately contains no hard-coded certification decision.
Certified wording is available only with an explicit release-rendering switch
and only when the sole receipt certifies the candidate, every hash binding and
mesh decision passes, and the band/DOS external-validation fields and their
HDF5 arrays are present and passing.
The current harmonic data status is \textbf{\HarmonicDataStatus}; physical mesh
convergence is \textbf{\PhysicalConvergenceStatus}.

\section{Conclusion}

The implemented 1.0 architecture closes a useful and well-defined bulk layer:
harmonic phonon observables, validated intrinsic three-phonon relaxation, and
SMRTA thermal conductivity with explicit units and approximation metadata.
\ifreportcertified The exact archived silicon cases demonstrate internal and
external implementation equality, not a converged material prediction. This
scientific certification does not by itself record the separately required
human PDF visual-inspection record or create a release tag.\else The provisional
silicon numbers demonstrate internal and external implementation equality,
not a converged material prediction. Completion of the exact clean release
gate, including band/DOS evidence and final visual inspection, is still
required before version 1.0.0 may be tagged.\fi{} Anharmonic NEGF junction
physics remains a distinct version 2.0 objective.

\begin{thebibliography}{9}

\bibitem{togo2015}
A. Togo, L. Chaput, and I. Tanaka,
``Distributions of phonon lifetimes in Brillouin zones,''
\emph{Physical Review B} \textbf{91}, 094306 (2015).

\bibitem{togo2023}
A. Togo, L. Chaput, T. Tadano, and I. Tanaka,
``Implementation strategies in phonopy and phono3py,''
\emph{Journal of Physics: Condensed Matter} \textbf{35}, 353001 (2023).

\bibitem{ziman1960}
J. M. Ziman,
\emph{Electrons and Phonons} (Oxford University Press, 1960).

\bibitem{guo2020}
Y. Guo \emph{et al.},
``Thermal transport in silicon films from anharmonic atomistic Green's
functions,'' \emph{Physical Review B} \textbf{102}, 195412 (2020).

\end{thebibliography}

\end{document}
